\section{Return counts and scope}
\label{sec:returns}

The complete orbit lists also determine the rational-point fixed sets
at every ordinary time. Let \(N_d=N_d(a,b)\) be the number of distinct
orbits of exact length \(d\), for \(d\in\{1,2,3,4\}\), obtained from
the appropriate table after removing its indicated degeneracies.
Put \(N_d=0\) for other positive integers \(d\).

\begin{corollary}[Ordinary-time returns]
\label{cor:returns}
For every \(a,b\in\Z\) and \(n\ge1\),
\begin{equation}
 \#\Fix(H_{a,b}^{\,n},\Q)
 =N_1+2N_2\ind{2\mid n}+3N_3\ind{3\mid n}
                  +4N_4\ind{4\mid n}.
 \label{eq:return_count}
\end{equation}
The rational-point Artin--Mazur series is
\begin{equation}
 \zeta_{a,b,\Q}(t)
 :=\exp\left(\sum_{n\ge1}
       \frac{\#\Fix(H_{a,b}^{\,n},\Q)}{n}t^n\right)
 =\prod_{d=1}^{4}(1-t^d)^{-N_d}.
 \label{eq:zeta}
\end{equation}
The equality holds as a formal power-series identity, and as an analytic
identity for \(|t|<1\).
\end{corollary}
\begin{proof}
A point on an orbit of exact length \(d\) is fixed by the \(n\)-th
iterate exactly when \(d\mid n\). Each such orbit contributes \(d\)
points, proving~\eqref{eq:return_count}. For each \(d\),
\[
 \sum_{j\ge1}\frac{d}{dj}t^{dj}
 =\sum_{j\ge1}\frac{t^{dj}}{j}=-\log(1-t^d).
\]
There are only finitely many orbits. Substitution and exponentiation
give~\eqref{eq:zeta} without any exchange of infinite orbit sums.
\end{proof}

For the equality locus in Theorem~\ref{thm:classification}, one has
\(N_2=1\), \(N_3=2\), and \(N_1=N_4=0\), so
\[
 \zeta_{a,b,\Q}(t)=\frac{1}{(1-t^2)(1-t^3)^2}.
\]
This is a rational-point orbit identity. It does not count all
algebraic fixed points, fixed-scheme lengths, or local intersection
multiplicities.

The proof identifies the source of the uniform bound within this
specific family. Integrality gives a discrete coordinate lattice;
the maximum-coordinate estimate then leaves six symbols; and the
coefficient parity separates the recurrence into two rigid local
equations. The integer and half-integer branches share this mechanism,
but their exceptional parameters and orbit coincidences differ.

The assumptions are not removed by the elementary nature of the
argument. For nonintegral rational coefficients the p-adic maximum
proof no longer supplies integrality, while other leading coefficients
or the opposite determinant sign change the recurrence. We make no
classification claim in those settings. In particular, the fixed
Jacobian-one result neither proves nor refutes the finite-exception
Jacobian conjecture in~\cite[Section 11]{dehenon2024problems}.
The broader rational-parameter and number-field problems remain
outside the theorem.

Finally, the finite product~\eqref{eq:zeta} carries only the ordinary
return data of the rational periodic set. No identification with target
Euler factors, functional-equation root numbers, automorphic data,
zeros of an arithmetic function, or a Hilbert--P\'olya realization is
asserted. The classification and the sharp eight-point bound are the
conclusions established here.
